\section{Problem Formulation and Theoretical Analysis}
\label{sec:analysis}
% 本章先限定离线视觉动作学习的对象与任务评价条件，
% 再分析动作拟合与实际任务表现之间的差距。
% 后续以动作的任务后果为共同尺度，

This section studies conditions for improving task performance from limited demonstrations. 
Starting from the relationship between imitation error and task success, we examine how to preserve information needed to distinguish the task consequences of different actions. 

We first establish conditions under which action error bounds task-failure probability 
and explain why improved coverage in representation space alone does not guarantee better task performance. 
% 随后分析语义引导如何保留专家选择依据，并由动作重构误差约束模仿误差迁移。
We then examine how semantic guidance can preserve task-relevant information
and, under consistent expert selection, control imitation-error transfer through expert-action reconstruction.
\subsection{Visuomotor Imitation Learning}
\label{sec:task}

We study offline visuomotor imitation learning from a finite, fixed set of robot demonstrations,
for manipulation tasks whose completion can be assessed from visual history.
The policy generates control commands from visual and proprioceptive observations and available execution history.
Training is based on this fixed demonstration dataset, with demonstrated actions providing the primary imitation target, 
whereas deployment performance is measured by the probability of task completion over a finite horizon.

\subsubsection{System and Observations}

Consider a robotic system governed by a single decision policy that performs a given manipulation task
through the joint actions of its actuators.
The joint robot-environment state $s_t\in\mathcal S$ evolves under the command
$u_t\in\mathcal A$ and gives rise to observation $o_t$.
The observation $o_t=(I_t^{1:V},q_t)$ comprises multi-view images
$I_t^{1:V}=(I_t^1,\ldots,I_t^V)$ and proprioceptive information $q_t$,
where $V\ge1$ is the number of cameras and $I_t^v$ is the image from the $v$th view.
Over a finite horizon of $T$ control steps, this interaction is represented by the following partially observed controlled process:
\begin{align}
 s_0&\sim p_0,\nonumber\\
 s_{t+1}&\sim P(\cdot\mid s_t,u_t),\quad 0\le t<T,\nonumber\\
 o_t&\sim O(\cdot\mid s_t),\quad 0\le t\le T.
 \label{eq:interactionmodel}
\end{align}
Here, $\mathcal S$ is the state space, $\mathcal A\subseteq\R^{d_u}$ is the admissible command set
specified by the chosen robot control interface,
$d_u$ is the command dimension, $p_0$ is the initial-state distribution,
and $P$ and $O$ are the transition and observation kernels, respectively.

\subsubsection{Offline Demonstrations and Policy}
The policy learns from a fixed demonstration set $\mathcal D=\{\zeta_i\}_{i=1}^{N}$
for the same manipulation task. The $i$th demonstration trajectory $\zeta_i$ records
observations and demonstrated control commands:
\begin{equation}
 \zeta_i=(o_0^i,u_0^{\mathrm E,i},o_1^i,\ldots,
 u_{T_i-1}^{\mathrm E,i},o_{T_i}^i).
 \label{eq:demonstrations}
\end{equation}
Here, $N$ is the finite number of demonstrations, $T_i\ge1$ is the number of control steps
in the $i$th trajectory, and the superscript $\mathrm E$ denotes demonstrated commands.


The information available to the deployed policy $\pi$ is represented by the history $\mathcal H_t$.
This history includes observations $o_{0:t}$ obtained up to time $t$, previously executed commands $u_{0:t-1}$, 
and other runtime records before the current decision.
The policy defines a distribution over admissible control commands,
with $u_t\sim\pi(\cdot\mid\mathcal H_t)$.
The command history is empty at $t=0$.

\subsubsection{Learning and Task Objectives}
For the following analysis, we consider deterministic policies acting on the available history,
with $\pi_{\mathrm E}$ denoting the expert policy.
Let $\nu$ denote the population distribution underlying the demonstrations,
where $(t,h)\sim\nu$ specifies a control step and its available history.
We define the population imitation risk as the expected squared
Euclidean deviation from the expert's action:
\begin{equation}
 R_\nu(\pi)
 :=\E_{(t,h)\sim\nu}\norm{\pi(h)-\pi_{\mathrm E}(h)}^2.
 \label{eq:populationimitationrisk}
\end{equation}
% 训练通过有限示范上的经验模仿损失来追求降低总体风险，具体损失在方法部分定义。
Training seeks to reduce this risk through an empirical imitation loss
on the finite demonstrations; the specific training losses are defined in Section~\ref{sec:method}.

The deployment objective is to complete the task within the finite horizon.
The measurable task-completion function $\Phi:\mathcal S^{T+1}\to\{0,1\}$ equals $1$
for state trajectories satisfying the task's terminal or temporal requirements and $0$ otherwise.
This defines the completion indicator $\Psi$ and the task-completion probability $J(\pi)$ of the deployed policy $\pi$:
\begin{equation}
 \Psi=\Phi(s_{0:T}),\qquad
 J(\pi)=\Prob_\pi(\Psi=1).
 \label{eq:visualtask}
\end{equation}
Here, $\Prob_\pi$ is the trajectory probability law induced jointly by the system model
and the complete deployed policy $\pi$.
% 学习旨在降低模仿风险，部署则追求更高任务成功率；后续分析研究二者何时能够联系起来。
Learning aims to reduce $R_\nu(\pi)$, whereas deployment aims to increase $J(\pi)$.
A lower imitation risk does not by itself guarantee a higher task-completion probability;
the following analysis examines conditions relating these objectives.


\subsubsection{Visual Assessability}
Under a fixed camera configuration and evaluation horizon $T$, the evaluation domain
$\Omega_{\mathrm{eval}}\subseteq\mathcal S^{T+1}\times\mathcal I^{T+1}$
consists of prespecified pairs of physical trajectories and visual histories
consistent with the system model and visibility conditions,
where $\mathcal I$ is the space of multi-view image tuples $I_t^{1:V}$.

\begin{assumption}[Visual assessability]
\label{ass:visualassessability}
% 存在可测的视觉判别函数 $\mathcal V:\mathcal I^{T+1}\to\{0,1\}$，
% 使任意 $(s_{0:T},I_{0:T}^{1:V})\in\Omega_{\mathrm{eval}}$ 均满足
There exists a measurable visual decision function $\mathcal V:\mathcal I^{T+1}\to\{0,1\}$
such that, for every $(s_{0:T},I_{0:T}^{1:V})\in\Omega_{\mathrm{eval}}$,
\begin{equation}
 \Phi(s_{0:T})=\mathcal V(I_{0:T}^{1:V}).
 \label{eq:visassumption}
\end{equation}
\end{assumption}

% 该假设只要求评价域内的任务完成可由视觉判断，不要求完整状态可观测。
Assumption~\ref{ass:visualassessability} specifies visual identifiability of completion within the evaluation domain,
without requiring full-state observability.
% 后续分析中的策略执行及反事实动作比较均限于物理轨迹与视觉历史几乎必然落在评价域内的过程。
The subsequent analysis considers policy executions and counterfactual action comparisons
for which $(s_{0:T},I_{0:T}^{1:V})\in\Omega_{\mathrm{eval}}$ almost surely.

\subsection{From Imitation Error to Task Performance}
\label{sec:actionanalysis}
Whether action error bounds task failure depends on where deviations occur and their subsequent task consequences.
We establish this connection and show that transferring imitation risk from demonstrations to deployment
depends jointly on representation-level coverage and conditional action-error discrepancy.

\subsubsection{Task-Performance Difference Between Policies}
\label{sec:actionconsequences}
We compare a policy $\pi$ with a reference policy $\pi_0$
through their task-failure probabilities and express the
difference as a sum of expected single-step action consequences.
Each consequence is defined by comparing alternative actions
at the same history while following $\pi_0$ afterward
\cite{ross2014aggrevate,cheng2018value}.
This decomposition provides the basis for relating imitation
error to task performance.

We consider deterministic policies acting on the available history.
Given a history $\mathcal H_t=h$ and a candidate action
$a\in\mathcal A$, define the action-value function as the
probability of task failure ($\Psi=0$):
\begin{equation}
 Q_t^{\pi_0}(h,a)
 :=\Prob^{a,\pi_0}\!\left(\Psi=0\mid\mathcal H_t=h\right).
 \label{eq:historyactionvalue}
\end{equation}
Here, $\Prob^{a,\pi_0}$ denotes the probability law obtained by
taking action $a$ at time $t$ and following $\pi_0$ afterward.

When the current action is also selected by $\pi_0$, this reduces
to the value function of $\pi_0$, $V_t^{\pi_0}(h) :=Q_t^{\pi_0}\bigl(h,\pi_0(h)\bigr).$
The effect of choosing $a$ instead of $\pi_0(h)$ is therefore
measured by
\begin{equation}
 A_t^{\pi_0}(h,a)
 :=Q_t^{\pi_0}(h,a)-V_t^{\pi_0}(h).
 \label{eq:actionconsequence}
\end{equation}
Thus, $A_t^{\pi_0}(h,a)>0$ indicates an increase in failure probability, whereas $A_t^{\pi_0}(h,a)<0$ indicates a decrease.

To relate these action-level consequences to differences in
task performance between policies, we require the following
common comparison conditions.
\begin{assumption}[Common task and environment]
\label{ass:commonprocess}
The compared policies share the same initial distribution and
task-completion criterion. Given a history $h$ and an action $a$,
the next available history follows a conditional distribution
that does not otherwise depend on the policy.
The reference policy $\pi_0$ is defined on all histories
involved in the comparison.
\end{assumption} 

Let $F(\pi)=1-J(\pi)$ denote the task-failure probability of policy $\pi$.
To evaluate $\pi$ relative to $\pi_0$, we set $a=\pi(h)$
and aggregate the resulting action consequences along trajectories generated by $\pi$.

Applying the finite-horizon performance-difference identity~\cite{kakade2002,ross2014aggrevate,cheng2018value}
to these one-step action effects gives
\begin{equation}
 F(\pi)-F(\pi_0)
 =\E_\pi\!\left[
 \sum_{t=0}^{T-1}A_t^{\pi_0}(\mathcal H_t,\pi(\mathcal H_t))
 \right].
\label{eq:pd}
\end{equation}
Here, $\E_\pi$ denotes expectation over trajectories generated by $\pi$.
Equation~\eqref{eq:pd} evaluates action effects at the situations actually visited by the evaluated policy $\pi$.
Its proof is given in Appendix~\ref{app:pdidentity}.
This relation allows task-performance differences to be
bounded through the consequences of action deviations.


\subsubsection{Task-Failure Bounds from Imitation Risk}
\label{sec:actionerrorbound}
Taking the expert policy as the reference,
we show that the imitation risk $R_\nu(\pi)$ bounds task-failure probability $F(\pi)$
when demonstrations sufficiently cover deployment histories
and action deviations have bounded task consequences.

Imitation risk is evaluated under the demonstration distribution $\nu$,
but deployment histories may follow a different distribution because actions depend on $\pi$.
Let $\mu_\pi^t$ denote the distribution of $(t,\mathcal H_t)$ at control step $t$ under $\pi$.
The distribution $\bar\mu_\pi=T^{-1}\sum_{t=0}^{T-1}\mu_\pi^t$ corresponds to choosing
a control step uniformly and observing its history.

% 对待评价策略，假设部署分布由示范分布以有限常数控制。
\begin{assumption}[Demonstration coverage]
\label{ass:democoverage}

For the evaluated policy $\pi$, there exists a uniform constant
$1\le C<\infty$ such that
\[
\bar\mu_\pi(\mathcal B)\le C\,\nu(\mathcal B)
\]
for every measurable set $\mathcal B$ of pairs $(t,h)$.
\end{assumption}
This condition excludes deployment situations absent from the reference distribution and limits how much their probabilities can be underestimated.

\begin{assumption}[Action sensitivity]
\label{ass:actionsensitivity}
Let $A_t^{\mathrm E}$ denote the action effect with the expert
policy as the reference.
For a fixed action scale, there exists a uniform constant
$0\le L<\infty$ such that
\[
|A_t^{\mathrm E}(h,a)|
\le L\norm{a-\pi_{\mathrm E}(h)}
\]
for all relevant $(t,h,a)$.
\end{assumption}
This assumption allows action deviations to be translated into bounds
on changes in task-failure probability.

Following the connection between deployment imitation error
and task cost~\cite{ross2011}, Assumptions~\ref{ass:democoverage}
and~\ref{ass:actionsensitivity}, together with the
performance-difference identity, yield the following bound (Appendix~\ref{app:actionriskbound}):
\begin{equation}
 F(\pi)\le F(\pi_{\mathrm E})+TL\sqrt{C R_\nu(\pi)}.
 \label{eq:successbound}
\end{equation}

Here, $C$ accounts for the shift from demonstration to deployment
distributions, while $L$ measures the sensitivity of task outcomes
to action deviations.
With other quantities fixed, a smaller $R_\nu(\pi)$ yields
a tighter upper bound on task-failure probability.

This bound identifies coverage and action sensitivity as key conditions
for translating imitation accuracy into task-performance guarantees.
We next examine the limitations of complete-history coverage
and the role of representations in risk transfer.

\subsubsection{Imitation-Risk Transfer Through Representations}
\label{sec:representationcoverage}
To address the limitations of complete-history coverage, we analyze
imitation-risk transfer in representation space.
We derive a bound showing that compression can relax coverage
requirements, while differences in conditional imitation error must also be accounted for.

Assumption~\ref{ass:democoverage} can fail even after an arbitrarily small action deviation.
For $T\ge2$, suppose the expert always executes $u_0=0$, whereas the learned policy executes
an admissible $u_0=\varepsilon\ne0$.
Histories recording $u_0=\varepsilon$ then occur under the learned policy
but have zero probability under the expert demonstration distribution.
No finite coverage constant $C$ can satisfy the assumption.

% 表示可合并上述历史差异，但也可能丢失区分动作后果所需的信息。
A representation can merge such distinctions, potentially relaxing coverage requirements,
but may also discard information needed to distinguish action consequences.
This tension between distributional alignment and information loss under a noninvertible representation
also appears in domain-adaptation analyses~\cite{johansson2019support,zhao2019invariant}.

% 固定策略、专家和同一个保留时间索引的可测因果编码器Z；mu_Z和nu_Z是部署与参考分布诱导的表示分布。
For a fixed policy, expert, and common measurable causal encoder $Z=\phi(t,\mathcal H_t)$ retaining the time index, let $\mu_Z$ and $\nu_Z$ be the representation distributions induced by $\bar\mu_\pi$ and $\nu$.
% 用表示层面的覆盖替代完整历史覆盖，不再要求假设3成立。
In place of Assumption~\ref{ass:democoverage}, suppose
$\mu_Z(\mathcal B)\le C_\phi\nu_Z(\mathcal B)$ for every measurable representation event $\mathcal B$, with $1\le C_\phi<\infty$.
% 同一表示下，两种分布仍可能包含不同历史；以D_phi衡量其条件平均动作平方误差的差异。
Even at the same representation, the two distributions may contain different histories.
Let $e(h)=\norm{\pi(h)-\pi_{\mathrm E}(h)}$ and define their conditional-error discrepancy as $D_\phi=\E_{z\sim\mu_Z}|m_\mu(z)-m_\nu(z)|$,
where $m_\mu(z)=\E_{\bar\mu_\pi}[e(h)^2\mid Z=z]$ and $m_\nu(z)=\E_\nu[e(h)^2\mid Z=z]$.
% 风险迁移只需要表示覆盖、平方可积和可测条件均值，不需要动作敏感性假设。
With square-integrable action error under both distributions and measurable versions of these conditional means,
\begin{equation}
 \E_{\bar\mu_\pi}e(h)^2
 \le C_\phi R_\nu(\pi)+D_\phi.
 \label{eq:representationrisktransfer-main}
\end{equation}
% 将动作风险转换成任务失败概率界时，仍需动作敏感性假设。
If Assumption~\ref{ass:actionsensitivity} also holds, the performance-difference identity gives
\begin{equation}
 F(\pi)\le F(\pi_{\mathrm E})
       +TL\sqrt{C_\phi R_\nu(\pi)+D_\phi}.
 \label{eq:representationcoveragebound}
\end{equation}

Appendix~\ref{app:representationcoverage} provides the proof.
% 固定模仿风险与动作敏感性时，表示对该界的影响体现为覆盖因子与条件误差差异。
With imitation risk and action sensitivity fixed,~\eqref{eq:representationcoveragebound}
shows that the effect of a representation on the bound depends jointly on
the coverage factor $C_\phi$ and the conditional error discrepancy $D_\phi$.
% 因此，一条收紧上界的途径是降低覆盖因子，同时控制条件误差差异，使组合项减小。
This motivates seeking representations that reduce $C_\phi$ while controlling $D_\phi$,
so that their combined contribution $C_\phi R_\nu(\pi)+D_\phi$ decreases.

\subsection{Task-Relevant Information Preservation}
\label{sec:representationanalysis}
% 承接B3：表示既要改善覆盖，又要保留解释专家动作所需的信息。
We seek conditions under which a representation can improve coverage while controlling the conditional imitation-error discrepancy.
% 语义用于引导信息保留，专家动作重构将这一设计思路连接到误差迁移界。
We examine semantic guidance as a way to retain task-relevant information, then connect the accuracy of that information to expert-action reconstruction and a bound on $D_\phi$.

% \subsubsection{语义引导的压缩与覆盖}
\subsubsection{Semantic-Guided Compression and Coverage}
% 压缩的目标是合并不影响动作选择的外观差异，而不是无条件追求分布对齐。
Semantic guidance can encourage an encoder to merge appearance differences while retaining object-level structure relevant to action selection.
% 这是表示设计动机，不是语义输入自动改善覆盖的结论。
This motivates a representation design, not an automatic coverage guarantee.
% 沿用B3，覆盖条件在表示空间的事件上成立，压缩后可能不再区分部分历史。
As in Section~\ref{sec:representationcoverage}, compression can relax coverage requirements by making some history distinctions irrelevant to representation events.
% 将确定性分割图附加在完整RGB上并不改变相对RGB的最优覆盖常数。
Appending a deterministic semantic map while retaining the original RGB input does not itself change the best coverage constant relative to RGB, since projection recovers that input.
% 因此分析对象是联合编码后的决策表示，而非输入通道数量。
The relevant object is the learned decision representation $Z=\phi(t,\mathcal H_t)$, not the number of input modalities.

% 分布合并若同时丢失专家动作选择依据，条件误差差异仍可能增大。
Improved coverage alone does not ensure a tighter bound; the error associated with merging histories must also be accounted for.
% 用共同专家读出重构动作是控制D的一个充分途径，不是必要条件。
A sufficient route is to retain information from which a common rule can reconstruct expert actions under both demonstration and deployment distributions.
% 这一行为相关的信息保留视角与既有表示抽象研究一致，但下文使用专家动作重构给出具体条件。
This emphasis on decision-relevant information is consistent with information-state and behaviorally grounded representation analyses~\cite{subramanian2022ais,zhang2021dbc}.
% 下文以对象属性和空间关系构造可检验的信息保留条件。
We next formulate this requirement using object attributes and spatial relationships.

% \subsubsection{任务信息与专家动作重构}
\subsubsection{Task Information and Expert-Action Reconstruction}
% 任务相关并不自动等于足以解释专家选择，需要明确近似条件。
We examine when retained task information is sufficient to approximate expert choices.
% 同时考虑信息本身的解释能力与从表示中估计这些信息的误差。
The reconstruction error depends on both what the selected information fails to explain and how accurately it is estimated.

% 固定被评价策略及两种分布，使用相同编码器和策略读出。
Fix the evaluated policy $\pi$ and set $\mu=\bar\mu_\pi$, with $\nu$ the demonstration distribution.
Write $\pi(h)=p(Z)$ for the same measurable policy readout under both distributions.
% G向量描述对象属性与空间关系；语义掩码可以提供相应监督量。
We define a task descriptor vector $G\in\R^{d_G}$ whose components describe task-relevant object attributes and spatial relationships.
Semantic masks provide a way to construct such quantities, including visibility and relative position.
% 上下文预先指定，而不是将完整历史不加限制地保留。
Let $\chi$ denote task-specified proprioceptive and execution context, selected in advance rather than identified with the unrestricted history.
% 读出估计和上下文均应来自策略实际使用的表示。
Suppose $Z$ retains $\chi$ and supports an estimate $\widehat G\in\R^{d_G}$ of $G$.
For $\rho\in\{\mu,\nu\}$, define $\delta_{G,\rho}=\E_\rho\norm{G-\widehat G}$.
% 两分布上均须定义这些量，包括缺失信息的明确约定。
These quantities must be defined under both distributions, with an explicit convention for unavailable components.

% \begin{assumption}[描述量条件下的专家选择]
\begin{assumption}[Descriptor-conditioned expert selection]
\label{ass:expertselection}
% 共同的专家近似规则在两种分布下使用同一组描述量和上下文。
There exists a common measurable rule $g_{\mathrm E}$ mapping descriptors and context to $\mathcal A$ such that, for $\rho\in\{\mu,\nu\}$,
\begin{equation}
 \E_\rho\norm{\pi_{\mathrm E}(h)-g_{\mathrm E}(G,\chi)}
 \le\epsilon_{\mathrm{sel},\rho}.
 \label{eq:expertselectionapprox}
\end{equation}
% 描述量估计误差对动作重构的影响由共同有限敏感性控制。
For a finite constant $K_{\mathrm E}\ge0$,
\begin{equation}
 \norm{g_{\mathrm E}(G,\chi)-g_{\mathrm E}(\widehat G,\chi)}
 \le K_{\mathrm E}\norm{G-\widehat G}
 \label{eq:expertselectionsensitivity}
\end{equation}
% 近似误差及估计误差有限，两种分布下几乎处处成立。
holds almost surely under both distributions, with finite $\epsilon_{\mathrm{sel},\rho}$ and $\delta_{G,\rho}$.
\end{assumption}
% 解释信息选择与信息估计各自承担的误差。
Here, $\epsilon_{\mathrm{sel},\rho}$ measures expert choices unexplained by $(G,\chi)$, while $K_{\mathrm E}$ controls sensitivity to estimation error.
% 任务有用性不自动保证专家选择一致性。
Task relevance alone does not imply this condition: equally successful actions may reflect different expert choices.

% 由表示读出的共同专家近似把描述量误差转为动作重构误差。
Define $\widehat a_{\mathrm E}(Z)=g_{\mathrm E}(\widehat G,\chi)$.
The triangle inequality yields
\begin{equation}
 \E_\rho\norm{\pi_{\mathrm E}(h)-\widehat a_{\mathrm E}(Z)}
 \le b_\rho
 :=\epsilon_{\mathrm{sel},\rho}+K_{\mathrm E}\delta_{G,\rho}.
 \label{eq:expertreconstruction}
\end{equation}
% 证明见附录；读出监督直接针对其中的描述量估计项。
Appendix~\ref{app:expertreconstruction} gives the derivation.
Supervising descriptor readouts targets the estimation term, while the choice of descriptors and context determines how well expert actions can be explained.

% \subsubsection{条件模仿误差差异的约束}
\subsubsection{Bounds on Conditional Imitation-Error Discrepancy}
\label{sec:relationtransfer}
% 共同的专家重构规则将信息保留直接连接到B3的D项。
Expert-action reconstruction connects information preservation to the discrepancy in Section~\ref{sec:representationcoverage}.
% 策略与专家近似都只依赖Z，两者的平方距离因而是跨分布相同的函数。
Because both $\pi(h)=p(Z)$ and $\widehat a_{\mathrm E}(Z)$ depend on the same representation, their squared difference is a common function under $\mu$ and $\nu$.
% 实际条件模仿误差偏离这一共同函数的程度由专家重构误差控制。
Expert reconstruction error bounds how far the actual conditional imitation errors depart from this function.

% 固定动作尺度，假设动作集合直径有限。
Assume a finite action-set diameter $B_{\mathcal A}=\sup_{a,a'\in\mathcal A}\norm{a-a'}$ under a fixed action scale.
% \begin{proposition}[通过专家动作重构约束条件误差迁移]
\begin{proposition}[Conditional-error transfer through expert reconstruction]
\label{prop:relationtransfer}
% 沿用专家选择与表示覆盖条件，重构误差上界由前式定义。
Under Assumption~\ref{ass:expertselection} and the representation-level coverage condition in Section~\ref{sec:representationcoverage}, with $b_\rho$ defined in~\eqref{eq:expertreconstruction},
\begin{equation}
 D_\phi\le 2B_{\mathcal A}(b_\mu+C_\phi b_\nu).
 \label{eq:relationdiscrepancybound}
\end{equation}
% 再以动作敏感性将模仿误差迁移转换为任务失败概率界。
If Assumption~\ref{ass:actionsensitivity} also holds,
\begin{equation}
 F(\pi)\le F(\pi_{\mathrm E})
 +TL\sqrt{C_\phi R_\nu(\pi)+2B_{\mathcal A}(b_\mu+C_\phi b_\nu)}.
 \label{eq:relationtransferfailure}
\end{equation}
\end{proposition}
% 证明对应附录，不再另行展开理想例子。
The proof is in Appendix~\ref{app:relationtransfer}.
% 回到B3的设计目标，同时限制覆盖加权误差和专家重构误差。
The design objective is therefore to improve coverage while preserving information needed to reconstruct expert actions.
% 比较任务界时固定外部条件，综合项下降才收紧上界。
For a fixed expert baseline, horizon, and action sensitivity, reducing the combined imitation and reconstruction terms tightens the task-failure bound.

% 提前提交双栏方法图，使其显示在方法正文附近的前一页。
\begin{figure*}[t]
\centering
\begin{tikzpicture}[>=Stealth,box/.style={draw,rounded corners,align=center,minimum height=1cm,font=\small}]
\node[box,text width=2.4cm] (sup) at (0,0) {Quality-aware\\semantic supervision};
\node[box,text width=3cm] (vis) at (4,1) {RGB, soft semantics,\\and proprioception};
\node[box,text width=3cm] (query) at (4,-1) {Supervised task-\\descriptor queries};
\node[box,text width=3cm] (enc) at (8,0) {Joint representation\\and ACT decoder};
\node[box,text width=2cm] (act) at (11.4,0) {Action-chunk\\prediction};
\draw[->] (sup) -- (vis);
\draw[->,dashed] (sup) -- (query);
\draw[->] (vis) -- (enc);
\draw[->] (query) -- (enc);
\draw[->] (enc) -- (act);
\end{tikzpicture}
\caption{Method overview. Quality-aware iterative annotation provides a frozen segmenter and quality-weighted descriptor labels. RGB, soft semantics, proprioception, and supervised queries jointly support action-chunk prediction.}
\label{fig:architecture}
\end{figure*}

